\documentclass[a4paper,11pt]{article}
\usepackage{fontspec}
\setmainfont{Latin Modern Roman}
\usepackage[ngerman]{babel}
\usepackage{amsmath,amssymb,amsthm,mathtools}
\usepackage{booktabs,tabularx,array}
\usepackage[a4paper,top=24mm,bottom=25mm,left=27mm,right=27mm]{geometry}
\usepackage{microtype}
\usepackage{needspace}
\usepackage[hidelinks]{hyperref}
\newcommand{\N}{\mathbb N}
\newcommand{\R}{\mathbb R}
\newcommand{\D}[1]{(\Delta^{#1}a)_0}
\newcommand{\GuideLink}{\href{../main-results/_B21-main-results.pdf\#formulaThm.69}{Theorem~21H.5.4.1}}
\setlength{\parindent}{0pt}
\setlength{\parskip}{.5em plus .1em minus .1em}
\setlength{\emergencystretch}{1.5em}
\allowdisplaybreaks[1]
\clubpenalty=10000
\widowpenalty=10000
\hypersetup{hypertexnames=false,pdftitle={Band 21 - Newtondarstellung und endliche Differenzen - Lesefassung},
 pdfauthor={Martin Kunze},pdfsubject={Lesefassung zu Theorem 21H.5.4.1 mit Prosabeweis, Beispielen und historischen Originalquellen}}
\title{Was die Differenzen einer Folge verraten\\[.5em]
\large Newtondarstellung und Abbruch der Differenzen\\
Lesefassung zu Theorem 21H.5.4.1 im Leitfaden für Folgen}
\author{Martin Kunze}
\date{}
\begin{document}
\maketitle
\hypertarget{newton.reading}{}
Die Quadratzahlen $0,1,4,9,16,\ldots$ wachsen nicht mit gleichbleibenden
Schritten. Ihre Zuwächse $1,3,5,7,\ldots$ tun es aber: Sie nehmen
jedes Mal um $2$ zu. Nach noch einer Differenz bleibt nur $0$.
Darin steckt ein allgemeines Prinzip. Jede Folge lässt sich aus ihren
Anfangsdifferenzen zurückgewinnen; endet die Differenzenbildung überall
nach derselben Anzahl von Schritten, so ist die Folge polynomial.

Diese Lesefassung erläutert und beweist \GuideLink{} vollständig.
Sie führt vom Zahlenbeispiel zu den tragenden Identitäten und bündelt sie
am Ende für die \href{_B21-newton-proofs.pdf\#newton.proofs}{separaten Beweistabellen}. Die Kennzeichen (N), (B) und
so weiter sind lokale Gleichungsmarken dieser Lesefassung; die
Satznummer im Leitfaden bleibt maßgeblich.

\section{Vom Zuwachs zur Differenzentafel}

Wir verwenden wie der Leitfaden $\N=\{0,1,2,\ldots\}$ und reelle Folgen
$a\colon\N\to\R$. Die \emph{Vorwärtsdifferenz}
\[
 (\Delta a)_n=a_{n+1}-a_n
\]
misst den Zuwachs beim Schritt von $n$ zu $n+1$. Ihre eigene Differenz
misst die Änderung dieses Zuwachses. Rekursiv setzen wir
\[
 (\Delta^0a)_n=a_n,\qquad
 (\Delta^{r+1}a)_n=(\Delta^ra)_{n+1}-(\Delta^ra)_n.
\]
In einer Differenzentafel steht jeder neue Eintrag für den rechten
minus den linken Nachbarn der vorigen Zeile. Für $a_n=n^2$ ergibt sich:
\begin{center}
\begin{tabular}{@{}lrrrrrr@{}}
\toprule
Index $n$ & $0$ & $1$ & $2$ & $3$ & $4$ & $5$\\
\midrule
$a_n$              & $\mathbf 0$ & $1$ & $4$ & $9$ & $16$ & $25$\\
$(\Delta a)_n$     & $\mathbf 1$ & $3$ & $5$ & $7$ & $9$ &\\
$(\Delta^2 a)_n$   & $\mathbf 2$ & $2$ & $2$ & $2$ & &\\
$(\Delta^3 a)_n$   & $\mathbf 0$ & $0$ & $0$ & & &\\
\bottomrule
\end{tabular}
\end{center}
Die fett gedruckten Anfangswerte $0,1,2,0$ werden die Koeffizienten
der Newtondarstellung. Leere Felder bedeuten nur, dass die sechs
angezeigten Ausgangswerte dort keine Differenz mehr liefern.
Für die \emph{ganze} Quadratfolge gilt dagegen
\[
 \Delta(n^2)=2n+1,\qquad \Delta^2(n^2)=2,\qquad \Delta^3(n^2)=0.
\]
Mit $\Delta(n^2)$ ist hier die Folge der Differenzen gemeint.
Eine endliche Tafel allein beweist noch kein Gesetz für alle Indizes.

\section{Die genaue Aussage}

Setze $D_k=(\Delta^ka)_0$. Für natürliche Indizes gilt
$\binom nk=n!/(k!(n-k)!)$ bei $0\leq k\leq n$ und
$\binom nk=0$ bei $k>n$; dabei ist $0!=1$.

\textbf{Newtondarstellung.} Für jede reelle Folge gilt
\begin{equation}\tag{N}\label{newton:identity}
 \boxed{\displaystyle a_n=\sum_{k=0}^{n}\binom nk D_k}
 \qquad(n\in\N).
\end{equation}
Jeder einzelne Wert wird durch eine \emph{endliche} Summe dargestellt.
Es werden weder ein Grenzwert noch Konvergenz einer Reihe vorausgesetzt.

\textbf{Abbruchkriterium.} Für ein festes $d\in\N$ sind äquivalent:
\begin{enumerate}
\item $\Delta^{d+1}a=0$, das heißt $(\Delta^{d+1}a)_n=0$ für jedes $n$;
\item $a_n=c_0+\sum_{j=1}^d c_jn^j$ für alle $n$, mit reellen $c_j$;
\item $a_n=\sum_{k=0}^d\binom nkD_k$ für alle $n$.
\end{enumerate}
Die Potenzkoeffizienten $c_0,\ldots,c_d$ sind eindeutig.
Unter diesen Bedingungen legen $a_0,\ldots,a_d$ die ganze Folge fest.
Für $d=0$ heißt das genau: $\Delta a=0$ und $a_n=a_0$.
Die Nullfolge ist eingeschlossen; ihr muss dafür kein eigener exakter
Grad zugewiesen werden.

Bei (N) darf die Zahl der benötigten Differenzen mit $n$ wachsen.
Beim Abbruchkriterium reicht dagegen \emph{ein festes} $d$ für alle $n$.
Das ist die entscheidende zusätzliche Aussage.

\section{Was bringt die Newtondarstellung beim Rechnen?}
\label{newton:practical}

Die Newtondarstellung ist sowohl ein Rechenverfahren als auch ein
Werkzeug für Beweise. Ihr praktischer Nutzen zeigt sich besonders,
wenn eine Folge durch \emph{Werte oder Differenzengesetze} gegeben ist:
Aus den Werten gewinnt man unmittelbar eine Formel; aus den
Differenzen gewinnt man ein einfaches Verfahren zum Weiterrechnen.
Ob die Darstellung die Rechnung verkürzt, hängt also davon ab, was
bereits bekannt ist und was berechnet werden soll.

\subsection*{Aus vier Werten eine Formel gewinnen}

Angenommen, aus der Aufgabenstellung ist bereits bekannt, dass die
\emph{ganze} Folge polynomial vom Grad höchstens $3$ ist. Gegeben seien
nur die ersten vier Werte
\[
 a_0=2,\qquad a_1=5,\qquad a_2=12,\qquad a_3=25.
\]
Gesucht sind eine Formel und etwa der weit entfernte Wert $a_{100}$.
Sechs Subtraktionen liefern die Differenzentafel:
\begin{center}
\begin{tabular}{@{}lrrrr@{}}
\toprule
Index $n$ & $0$ & $1$ & $2$ & $3$\\
\midrule
$a_n$            & $\mathbf2$ & $5$ & $12$ & $25$\\
$(\Delta a)_n$   & $\mathbf3$ & $7$ & $13$ &\\
$(\Delta^2a)_n$  & $\mathbf4$ & $6$ & &\\
$(\Delta^3a)_n$  & $\mathbf2$ & & &\\
\bottomrule
\end{tabular}
\end{center}
Die Anfangsspalte gibt sofort $D_0=2$, $D_1=3$, $D_2=4$, $D_3=2$.
Wegen der vorausgesetzten Gradschranke sind alle höheren Differenzen
null. Daher gilt für jedes $n$
\[
 \boxed{\displaystyle a_n=2+3n+4\binom n2+2\binom n3.}
\]
Wir können direkt zum gewünschten Index springen:
\[
 a_{100}=2+3\cdot100+4\cdot4950+2\cdot161700=343502.
\]
Dazu mussten die Werte $a_4,\ldots,a_{99}$ nicht einzeln berechnet
werden. Der eigentliche Gewinn beim Aufstellen der Formel ist:
Wir lesen ihre Koeffizienten aus der Differenzentafel ab. Beim Ansatz
$a_n=c_0+c_1n+c_2n^2+c_3n^3$ müssten wir zunächst ein Gleichungssystem
für die vier Potenzkoeffizienten lösen.

Ausmultiplizieren ergibt hier auch
\[
 a_n=\frac13n^3+n^2+\frac53n+2.
\]
Ist diese Potenzform bereits bekannt, kann man einen einzelnen Wert
ebenfalls unmittelbar berechnen, etwa mit dem Hornerverfahren.
Die Newtonform ist also nicht grundsätzlich die schnellste Auswertung;
sie ist hier besonders leicht \emph{aus den gegebenen Werten zu gewinnen}.

\subsection*{Viele aufeinanderfolgende Werte: nur addieren}

Sollen nun $a_4,a_5,a_6,\ldots$ der Reihe nach erzeugt werden, braucht
man die Formel nicht jedes Mal neu auszuwerten. Schreibe
$u_n=(\Delta a)_n$ und $v_n=(\Delta^2a)_n$.
Die dritte Differenz ist im Beispiel überall $2$, also gilt
\[
 \begin{aligned}
 a_{n+1}&=a_n+u_n,\\
 u_{n+1}&=u_n+v_n,\\
 v_{n+1}&=v_n+2.
 \end{aligned}
\]
Alle rechten Seiten benutzen die Werte beim alten Index $n$.
Aus dem Anfangszustand $(a_0,u_0,v_0)=(2,3,4)$ erhält man so
\begin{center}
\begin{tabular}{@{}rrrr@{}}
\toprule
$n$ & $a_n$ & $u_n$ & $v_n$\\
\midrule
$0$ & $2$  & $3$  & $4$\\
$1$ & $5$  & $7$  & $6$\\
$2$ & $12$ & $13$ & $8$\\
$3$ & $25$ & $21$ & $10$\\
$4$ & $46$ & $31$ & $12$\\
$5$ & $77$ & $43$ & $14$\\
\bottomrule
\end{tabular}
\end{center}
Jede weitere Tabellenzeile kostet genau drei Additionen.
Allgemein reichen bei einer Folge vom Grad höchstens $d$ nach der
Vorbereitung $d$ Additionen pro Schritt; man führt die aktuellen
Differenzen bis zur konstanten $d$-ten Differenz mit.
Das ist ein konkreter Vorteil beim Erzeugen langer Wertetabellen.
Die Newtonkoeffizienten liefern genau den benötigten Anfangszustand.

\subsection*{Auch Summieren und Beweisen werden einfacher}

Die Binomialbasis passt zu Differenzen und Summen. Die später
bewiesene Summationsformel \eqref{newton:sum} liefert im selben Beispiel
unmittelbar
\[
 \sum_{n=0}^{N-1}a_n
 =2N+3\binom N2+4\binom N3+2\binom N4.
\]
Man erhöht in jedem Binomialterm nur den unteren Index um eins.
Das ergibt eine geschlossene Formel für beliebig viele Summanden,
ohne sie einzeln zu addieren oder zuerst Potenzsummen herzuleiten.

Für Beweise ist derselbe Aufbau nützlich: Differenzen senken den
Basisindex, Anfangsdifferenzen lesen Koeffizienten aus, und der globale
Abbruch kennzeichnet die polynomialen Folgen. So lassen sich
Rekonstruktion, Eindeutigkeit und Summenformeln mit denselben wenigen
Rechenregeln behandeln.

\subsection*{Was man aus endlich vielen Werten nicht gewinnt}

Die Gradannahme im Beispiel ist entscheidend. Vier Werte allein
beweisen nicht, dass die Folge kubisch ist oder dass alle späteren
Werte der gefundenen Formel folgen. Ohne diese Zusatzinformation
beschreibt die Formel zunächst nur das Interpolationspolynom durch
die vier gegebenen Werte. Eine Fortsetzung damit ist dann eine
Modellannahme, keine Folgerung aus den vier Zahlen.

Für eine beliebige Folge ist die allgemeine Identität (N) noch kein
Verfahren, unbekannte Werte vorherzusagen: Zur Berechnung von
$D_n=(\Delta^na)_0$ aus der Tafel braucht man bereits $a_n$.
Eine Vereinfachung entsteht erst, wenn die Differenzen zusätzliche
Struktur besitzen oder unabhängig bekannt sind. Und ist schon eine
kurze Formel wie $a_n=2^n$ gegeben, macht ihre Umschreibung als
$\sum_{k=0}^n\binom nk$ die Auswertung gewöhnlich aufwendiger.
Die passende Darstellung richtet sich nach der Aufgabe: Newton ist
besonders hilfreich beim Gewinnen einer Formel aus Werten, beim
Fortsetzen polynomialer Wertetabellen sowie beim Differenzieren und
Summieren von Folgen.

\section{Woher die Binomialkoeffizienten kommen}

Die Rekursionsgleichung der Differenzentafel lässt sich rückwärts lesen:
\begin{equation}\tag{R}\label{newton:reconstruction}
 (\Delta^ra)_{n+1}=(\Delta^ra)_n+(\Delta^{r+1}a)_n.
\end{equation}
Um eine Spalte nach rechts zu gelangen, addiert man also den Eintrag
und den unmittelbar darunterliegenden Eintrag der alten Spalte.
Wiederholung erzeugt das Pascal-Dreieck:
\[
 \begin{aligned}
 a_0&=D_0,\\
 a_1&=D_0+D_1,\\
 a_2&=D_0+2D_1+D_2,\\
 a_3&=D_0+3D_1+3D_2+D_3.
 \end{aligned}
\]
Der Term $D_k$ trägt nach $n$ Schritten so oft bei, wie man unter
$n$ Schritten die $k$ Schritte auswählen kann, die zur nächsttieferen
Differenz führen. Diese Anzahl ist $\binom nk$.

Die zugrunde liegende Rechenregel lautet
\begin{equation}\tag{P}\label{newton:pascal}
 \binom{n+1}{k}=\binom nk+\binom n{k-1}.
\end{equation}
Für die Randfälle setzen wir hier zusätzlich $\binom n{-1}=0$.
Für $1\leq k\leq n$ folgt (P) durch gemeinsames Nennerbilden aus der
Fakultätsformel; $k=0,n+1$ sind unmittelbar. Anschaulich wird eine
$k$-elementige Auswahl danach getrennt, ob sie das letzte Element
enthält. Mehr als diese Identität braucht der erste Beweis nicht.

\subsection*{Ein kurzer Beweis für jede Folge}

Wir beweisen (N) durch Induktion über $n$, \emph{gleichzeitig für alle
reellen Folgen}. Für $n=0$ ist sie $a_0=D_0$. Die Induktionsannahme
lautet: Für jede reelle Folge $b$ gilt
\[
 b_n=\sum_{k=0}^n\binom nk(\Delta^kb)_0.
\]
Wir wenden diese Aussage zweimal an. Für $b=a$ ergibt sie die
bekannte Summe $a_n=\sum_{k=0}^n\binom nkD_k$.
Für die zweite Anwendung setzen wir ausdrücklich $b=\Delta a$.
Die Anfangsdifferenzen dieser neuen Folge sind
\[
 b_0=D_1,\qquad (\Delta b)_0=D_2,\qquad
 (\Delta^kb)_0=(\Delta^k(\Delta a))_0
               =(\Delta^{k+1}a)_0=D_{k+1}.
\]
Denn $b$ ist bereits die erste Differenzenfolge von $a$; jede weitere
Differenz erhöht deren Ordnung um eins. Die Induktionsannahme liefert
also gerade den zweiten Summanden der folgenden Rechnung:
\[
 \boxed{\displaystyle
 (\Delta a)_n=b_n
 =\sum_{k=0}^n\binom nk(\Delta^kb)_0
 =\sum_{k=0}^n\binom nkD_{k+1}.}
\]
Der Index $n$ und die Faktoren $\binom nk$ bleiben gleich: Wir berechnen
weiterhin das $n$-te Glied, jetzt aber dasjenige der Folge $\Delta a$.
Nur ihre Anfangsdifferenzen beginnen bei $D_1$ statt bei $D_0$.
Zum Beispiel ergibt sich bei $n=2$
\[
 a_2=D_0+2D_1+D_2,\qquad
 (\Delta a)_2=D_1+2D_2+D_3.
\]
Ihre Summe ist $a_3=D_0+3D_1+3D_2+D_3$.

Um die beiden Summen allgemein zusammenzufassen, benennen wir in der
zweiten den Index um: $j=k+1$. Dann wird aus $D_{k+1}$ ein $D_j$,
aus $\binom nk$ ein $\binom n{j-1}$ und aus $0\leq k\leq n$ der
Bereich $1\leq j\leq n+1$. Damit gilt
\[
 \sum_{k=0}^n\binom nkD_{k+1}
 =\sum_{j=1}^{n+1}\binom n{j-1}D_j.
\]
Der Buchstabe des Summationsindex ist beliebig; rechts schreiben wir
anschließend wieder $k$. Nun lautet die ganze Gleichungskette
\[
\begin{aligned}
 a_{n+1}
 &=a_n+(\Delta a)_n\\
 &=\sum_{k=0}^n\binom nkD_k+
   \sum_{k=0}^n\binom nkD_{k+1}\\
 &=\sum_{k=0}^{n+1}\binom nkD_k+
   \sum_{k=0}^{n+1}\binom n{k-1}D_k\\
 &=\sum_{k=0}^{n+1}
   \left(\binom nk+\binom n{k-1}\right)D_k\\
 &=\sum_{k=0}^{n+1}\binom{n+1}{k}D_k.
\end{aligned}
\]
In der dritten Zeile wurden nur Nullsummanden ergänzt:
$\binom n{n+1}D_{n+1}=0$ in der ersten und
$\binom n{-1}D_0=0$ in der zweiten Summe. Die letzte Zeile ist (P).
Damit ist die Induktion geschlossen. Die gleichzeitige Quantifizierung
über alle Folgen ist wesentlich: Eine Induktionsannahme nur für das
festgehaltene $a$ dürfte nicht ohne Weiteres für $\Delta a$ benutzt werden.

\subsection*{Dieselbe Idee als binomischer Lehrsatz}

Wer Operatoren kennt, kann den Beweis auf eine Zeile verdichten.
Die Verschiebung $E$ ist durch $(Ea)_n=a_{n+1}$ gegeben, die Identität
$I$ durch $Ia=a$. Es gilt $E=I+\Delta$. Weil $I$ und $\Delta$
kommutieren, liefert der binomische Lehrsatz
\[
 a_n=(E^na)_0=((I+\Delta)^na)_0
     =\sum_{k=0}^n\binom nkD_k.
\]
Der Induktionsbeweis oben enthält genau dieses Argument ohne eine
zusätzliche Theorie linearer Operatoren und eignet sich daher besser
als Grundlage der Beweistabellen.

\section{Warum gerade diese Bausteine passen}

Die gewöhnlichen Potenzen $1,x,x^2,\ldots$ sind für Differenzen
weniger bequem: Schon $\Delta(n^3)=3n^2+3n+1$ besteht aus mehreren
Termen. Zweckmäßiger sind die \emph{Binomialpolynome}
\begin{equation}\tag{B}\label{newton:basis}
 B_0(x)=1,\qquad
 B_k(x)=\binom xk:=\frac{x(x-1)\cdots(x-k+1)}{k!}\quad(k\geq1).
\end{equation}
Hier bezeichnet $\binom xk$ ein Polynom auch für reelles $x$.
An natürlichen Stellen stimmt es mit der vorigen Konvention überein;
für $k>n$ enthält das Produkt den Faktor $n-n=0$.
Setzen wir $B_{-1}=0$, so gilt für alle $k\geq0$
\begin{equation}\tag{S}\label{newton:lowering}
 B_k(x+1)-B_k(x)=B_{k-1}(x).
\end{equation}
Für $k\geq2$ zieht man aus der linken Seite
$x(x-1)\cdots(x-k+2)/k!$ heraus. Der verbleibende Faktor ist
$(x+1)-(x-k+1)=k$. Für $k=0,1$ liest man die Gleichung direkt ab.
Damit ist (S) auch unabhängig von einer Beschränkung auf natürliche
Stellen bewiesen.

\textbf{Differenzieren heißt hier nur: einen Index heruntergehen.}
Insbesondere verschwinden nach $k+1$ Differenzen alle Werte von $B_k$.
Außerdem ist $B_0(0)=1$ und $B_k(0)=0$ für $k\geq1$. Wiederholte
Anwendung von (S) ergibt für die Folgen $B_k(n)$
\begin{equation}\tag{K}\label{newton:coordinates}
 (\Delta^j B_k)_0=
 \begin{cases}1,&j=k,\\0,&j\ne k.\end{cases}
 \qquad(j,k\in\N).
\end{equation}
Für $j>k$ ist die Differenzenfolge bereits null.
Da $\Delta$ linear ist, folgt daraus sofort:
\[
 p(n)=\sum_{k=0}^d\beta_kB_k(n)
 \quad\Longrightarrow\quad (\Delta^jp)_0=\beta_j
 \quad(0\leq j\leq d).
\]
Die Anfangsdifferenz liest also genau \emph{einen} Koeffizienten aus.
Man muss dazu kein Gleichungssystem lösen.

\subsection*{Was an den Differenzentermen besonders ist}

Drei verschiedene Dinge sollten auseinandergehalten werden:
$(\Delta^ka)_n$ ist eine ganze Differenzenfolge, $D_k$ ihr Anfangswert
und $\binom nkD_k$ der zugehörige Summand der Newtondarstellung.
Der Anfangswert ist eine aus den Daten berechnete Koordinate;
das Binomialpolynom beschreibt, wie diese Koordinate zu späteren
Folgengliedern beiträgt.

Die ersten Koeffizienten lauten
\[
 D_0=a_0,\quad D_1=a_1-a_0,\quad
 D_2=a_2-2a_1+a_0,\quad D_3=a_3-3a_2+3a_1-a_0.
\]
Allgemein gilt die endliche Formel
\begin{equation}\tag{D}\label{newton:explicit}
 (\Delta^ka)_n=\sum_{j=0}^k(-1)^{k-j}\binom kj a_{n+j}.
\end{equation}
Ein Induktionsschritt besteht darin, die Formel an $n+1$ und an $n$
zu subtrahieren und mit (P) zusammenzufassen; für $k=0$ ist sie
unmittelbar. Alternativ ist sie die Entwicklung von $(E-I)^k$.
Die wechselnden Vorzeichen sorgen für die Auslöschung niedriger
Polynomgrade. Insbesondere ändert das Addieren einer polynomialen
Folge vom Grad kleiner als $k$ die $k$-te Differenz nicht.

Jedes $D_k$ benutzt nur $a_0,\ldots,a_k$. Nach (N) ist
\begin{equation}\tag{C}\label{newton:correction}
 D_m=a_m-\sum_{k=0}^{m-1}\binom mkD_k\quad(m\geq1).
\end{equation}
Es misst genau den Fehler, den die bisherigen Bausteine an der
nächsten Stelle $m$ noch machen. Die Korrektur $D_mB_m$ lässt alle
früheren Stellen unverändert, weil $B_m(0)=\cdots=B_m(m-1)=0$, und
korrigiert die Stelle $m$ vollständig, weil $B_m(m)=1$.
Das erklärt die Darstellung auch als schrittweise Interpolation.

\section{Der kurze Beweis des Abbruchkriteriums}

\textbf{(1) $\Rightarrow$ (3).}
Aus $\Delta^{d+1}a=0$ folgt durch weitere Differenzenbildung
$\Delta^ka=0$ für jedes $k>d$. In (N) verschwinden daher alle
entsprechenden Summanden. Ist $n<d$, ergänzen wir Nullsummanden,
weil $\binom nk=0$ für $k>n$. So gilt die feste Summe aus (3) für
\emph{alle} $n$.

\textbf{(3) $\Rightarrow$ (2).}
Jedes $B_k(x)$ aus (B) ist ein Polynom vom Grad $k$ mit Leitkoeffizient
$1/k!$. Die endliche Summe $\sum_{k=0}^d D_kB_k(x)$ ist deshalb ein
Polynom vom Grad höchstens $d$. Ausmultiplizieren liefert (2).

\Needspace{10\baselineskip}
\textbf{(2) $\Rightarrow$ (1).}
Eine Differenz senkt den Grad eines nichtkonstanten Polynoms um genau
eins. Denn für $j\geq1$ ist nach dem binomischen Lehrsatz
\begin{equation}\tag{G}\label{newton:degree}
 (n+1)^j-n^j=\sum_{i=0}^{j-1}\binom ji n^i,
\end{equation}
mit höchstem Term $j n^{j-1}$. Konstanten werden zu null.
Linearität und $d+1$ Wiederholungen liefern daher $\Delta^{d+1}a=0$.
Das beweist die drei Äquivalenzen.

\textbf{Eindeutigkeit der Potenzkoeffizienten.}
Hätten zwei Polynome vom Grad höchstens $d$ dieselben Folgenwerte,
so wäre ihre Differenz überall null. Wäre $c_rx^r$ ihr höchster
nichtverschwindender Term, so wäre ihre $r$-te Differenz nach (G)
die von null verschiedene Konstante $r!c_r$. Das ist unmöglich.
Für $r=0$ ist bereits das konstante Differenzpolynom selbst ungleich
null. Also stimmen alle Koeffizienten überein.

\textbf{Bestimmung durch Anfangswerte.}
Die Werte $a_0,\ldots,a_d$ bestimmen durch die Differenzentafel
$D_0,\ldots,D_d$, und (3) bestimmt daraus jedes $a_n$.
Hier ist die Annahme $\Delta^{d+1}a=0$ für alle Indizes unverzichtbar.
Damit sind auch die Zusatzbehauptungen des Theorems bewiesen.

Ist die Folge nicht null und ihr exakter Grad $r$, so gilt sogar
\[
 \Delta^r a=r!c_r\ne0,\qquad \Delta^{r+1}a=0.
\]
Die letzte nichtverschwindende Differenzenzeile ist konstant.
Ihr Wert ist der höchste Newtonkoeffizient $D_r$; der höchste
Potenzkoeffizient ist $c_r=D_r/r!$.

\section{Beispiele, die unterschiedliche Seiten sichtbar machen}

\subsection*{1. Quadrate: dieselbe Folge in einer passenden Basis}

Die Anfangsspalte der ersten Tafel lautet $D_0=0$, $D_1=1$,
$D_2=2$, danach null. Deshalb ist
\[
 n^2=\binom n1+2\binom n2=n+n(n-1).
\]
Die zweite Differenz ist $2$, obwohl der Leitkoeffizient von $n^2$
gleich $1$ ist: Der Faktor $2!=2$ erklärt den Unterschied.
Für $n=5$ liest man $25=5+2\cdot10$ ab.

\subsection*{2. Kuben: die Koeffizienten aus der Anfangsspalte}

Für $a_n=n^3$ ergibt sich die Tafel
\begin{center}
\begin{tabular}{@{}lrrrrrr@{}}
\toprule
Index $n$ & $0$ & $1$ & $2$ & $3$ & $4$ & $5$\\
\midrule
$a_n$            & $\mathbf0$ & $1$ & $8$ & $27$ & $64$ & $125$\\
$(\Delta a)_n$   & $\mathbf1$ & $7$ & $19$ & $37$ & $61$ &\\
$(\Delta^2a)_n$  & $\mathbf6$ & $12$ & $18$ & $24$ & &\\
$(\Delta^3a)_n$  & $\mathbf6$ & $6$ & $6$ & & &\\
$(\Delta^4a)_n$  & $\mathbf0$ & $0$ & & & &\\
\bottomrule
\end{tabular}
\end{center}
Der globale Abbruch folgt hier aus dem bekannten Polynom $n^3$,
nicht nur aus den sichtbaren Nullen. Newton liefert
\[
 n^3=\binom n1+6\binom n2+6\binom n3.
\]
Zum Beispiel ist $64=4+6\cdot6+6\cdot4$.
Ausmultiplizieren würde die drei Newtonkoeffizienten wieder in
die einzige Potenz $n^3$ zusammenfassen; für die Differenzenbildung
ist gerade die zerlegte Form einfacher.

\subsection*{3. Ganzzahlige Werte trotz gebrochener Potenzkoeffizienten}

Die Dreieckszahlen $T_n=n(n+1)/2$ beginnen mit $0,1,3,6,10,\ldots$.
Ihre Anfangsdifferenzen sind $0,1,1,0,\ldots$, also
\[
 T_n=\binom n1+\binom n2=\frac12n^2+\frac12n.
\]
Die Newtonkoeffizienten sind ganzzahlig, die Potenzkoeffizienten nicht.
Allgemein hat ein reelles Polynom genau dann ganzzahlige Werte auf
$\N$, wenn alle seine Newtonkoeffizienten ganzzahlig sind:
Die eine Richtung folgt aus der Differenzentafel, die andere daraus,
dass $\binom nk$ ganzzahlig ist. Die Newtonbasis bewahrt somit eine
arithmetische Information, die in der Potenzbasis weniger sichtbar ist.

\subsection*{4. Eine geometrische Folge: Newton ohne Abbruch}

Für $a_n=2^n$ ist $\Delta a=a$. Somit gilt $D_k=1$ für jedes $k$,
und (N) wird zu
\[
 2^n=\sum_{k=0}^n\binom nk.
\]
Die Formel ist für jedes $n$ endlich und exakt, obwohl keine
Differenzenordnung die ganze Folge auslöscht. Die Folge ist daher
nicht polynomial. Allgemeiner gilt für $a_n=q^n$ bei $q\ne0$
\[
 (\Delta^ka)_n=(q-1)^kq^n,\qquad
 q^n=\sum_{k=0}^n\binom nk(q-1)^k.
\]
Hier erscheint der gewöhnliche binomische Lehrsatz unmittelbar wieder.

\subsection*{5. Eine Tafel kann über den späteren Verlauf täuschen}

Betrachte die beiden Folgen
\[
 a_n=n^2,\qquad b_n=n^2+\binom n6.
\]
Für $n=0,\ldots,5$ sind sie gleich; sämtliche aus diesen sechs
Werten berechenbaren Differenzen stimmen überein. Bei $n=6$ gilt
aber $a_6=36$, $b_6=37$. Für die zweite Folge ist
$\Delta^3b_n=\binom n3$ und insbesondere $\Delta^6b=1$.
Ihr Grad ist $6$, nicht $2$.

Der zusätzliche Summand ist unsichtbar, bis sein Index erreicht wird.
Allgemein verändert $\lambda\binom n{m+1}$ keinen der Werte
$a_0,\ldots,a_m$, wohl aber für $\lambda\ne0$ den nächsten Wert.
Ein endliches Anfangsstück bestimmt also seine Interpolation vom
Grad höchstens $m$, jedoch ohne Zusatzannahme nicht die ganze Folge.

\Needspace{18\baselineskip}
\section{Interpolation, Summation und der Vergleich mit Taylor}

Aus beliebigen Daten $a_0,\ldots,a_m$ entsteht das Polynom
\[
 P_m(x)=\sum_{k=0}^mD_kB_k(x).
\]
Für $0\leq n\leq m$ ist $P_m(n)=a_n$ nach (N).
Es ist das einzige Polynom vom Grad höchstens $m$ mit diesen Werten:
Bei der Differenz zweier solcher Polynome liefern die ersten
$m+1$ Nullwerte auch $D_0=\cdots=D_m=0$; die bereits bewiesene
Polynomfassung von (N) macht die Differenz an allen natürlichen Stellen null. Nach der bereits bewiesenen Eindeutigkeit der Potenzkoeffizienten ist sie daher das Nullpolynom.
Dies ist \emph{Interpolation}. Eine Behauptung über noch unbekannte
Folgenwerte wird daraus erst durch eine zusätzliche Voraussetzung.

\subsection*{Ein anderer Anfangspunkt und eine andere Schrittweite}

Für $s\in\N$ wendet man (N) auf die Restfolge $a_{s+n}$ an und erhält
\[
 a_{s+n}=\sum_{k=0}^n\binom nk(\Delta^ka)_s.
\]
Für Daten an gleichmäßig verteilten Stellen $x_0+nh$ mit $h\ne0$
setzt man $a_n=f(x_0+nh)$ und
$\Delta_h f(x)=f(x+h)-f(x)$. Das Interpolationspolynom lautet dann
\[
 P_m(x)=\sum_{k=0}^m
 \frac{\prod_{j=0}^{k-1}(x-x_0-jh)}{k!h^k}
 (\Delta_h^k f)(x_0),
\]
wobei das leere Produkt $1$ ist. Die Formel folgt durch die
Substitution $n=(x-x_0)/h$ in $P_m$; sie interpoliert die
$m+1$ gegebenen Gitterwerte. Für beliebiges $f$ wird damit keine
Identität an allen reellen Stellen behauptet.

\subsection*{Die diskrete Entsprechung zur Taylordarstellung}

Die Taylordarstellung entwickelt eine Funktion um eine Stelle $x_0$
in Potenzen von $x-x_0$. Für ein Polynom $p$ vom Grad höchstens $d$
ist sie die endliche, exakte Identität
\begin{equation}\tag{T}\label{newton:taylor}
 \boxed{\displaystyle
 p(x)=\sum_{k=0}^d p^{(k)}(x_0)\frac{(x-x_0)^k}{k!}.}
\end{equation}
Dabei bezeichnet $p^{(k)}$ die $k$-te Ableitung und $p^{(0)}=p$.
Warum stehen hier gerade diese Koeffizienten? Setze
$\varphi_k(x)=(x-x_0)^k/k!$. Diese normierten Potenzen erfüllen
\[
 \varphi_k'=\varphi_{k-1}\quad(k\geq1),\qquad
 \varphi_k^{(j)}(x_0)=
 \begin{cases}1,&j=k,\\0,&j\ne k.\end{cases}
\]
Jede Ableitung senkt den Basisindex um eins. Nach $j$ Ableitungen und
Einsetzen von $x_0$ bleibt nur der Baustein mit Index $j$ übrig.
Schreibt man $p=\sum_{k=0}^d \alpha_k\varphi_k$, so folgt deshalb sofort
$\alpha_j=p^{(j)}(x_0)$. Das beweist (T): Die verschobenen Potenzen bilden
eine Basis der Polynome vom Grad höchstens $d$, und das Ableiten liest
ihre Koeffizienten aus.

Bei Newton übernehmen die Binomialpolynome dieselbe Aufgabe:
\[
 \Delta B_k=B_{k-1},\qquad
 (\Delta^jB_k)_0=
 \begin{cases}1,&j=k,\\0,&j\ne k.\end{cases}
\]
Hier liest die $j$-te Anfangsdifferenz den Koeffizienten von $B_j$ aus.
Die beiden Darstellungen eines Polynoms bei Anfangspunkt $0$ sind also
\[
 \underbrace{p(x)=\sum_{k=0}^d p^{(k)}(0)\frac{x^k}{k!}}_{\text{Taylor: Ableitungen und normierte Potenzen}},
 \qquad
 \underbrace{p(x)=\sum_{k=0}^d(\Delta^kp)_0B_k(x)}_{\text{Newton: Differenzen und Binomialpolynome}}.
\]
Für Polynome gilt auch die zweite Gleichung für alle reellen $x$:
Die Gleichheit an allen natürlichen Stellen und die Eindeutigkeit der
Potenzkoeffizienten zeigen, dass beide Seiten dasselbe Polynom sind.

\textbf{Dasselbe Polynom, unterschiedliche Koeffizienten.}
Für $p(x)=x^2$ bei $x_0=0$ lauten die Anfangsableitungen
$(p(0),p'(0),p''(0))=(0,0,2)$, die Anfangsdifferenzen dagegen
$(D_0,D_1,D_2)=(0,1,2)$. Daher
\[
 \begin{aligned}
 \text{Taylor:}\quad x^2&=0\cdot1+0\cdot x+2\frac{x^2}{2!},\\
 \text{Newton:}\quad x^2&=0\cdot1+1\cdot x+2\frac{x(x-1)}{2!}.
 \end{aligned}
\]
Die Ableitung $p'(0)=0$ misst die momentane Steigung bei $0$.
Die Differenz $(\Delta p)_0=p(1)-p(0)=1$ misst den Zuwachs über einen
ganzen Schritt. Beide Koeffizienten erfüllen dieselbe Aufgabe in ihrer
jeweiligen Basis; sie müssen deshalb nicht denselben Wert haben.

\subsection*{Wie Taylor bei kleiner Schrittweite entsteht}

In der Newtonformel mit Schrittweite $h$ lassen sich die Faktoren so
anordnen:
\[
 P_{m,h}(x)=\sum_{k=0}^m
 \frac{(\Delta_h^kf)(x_0)}{h^k}\,
 \frac{\prod_{j=0}^{k-1}(x-x_0-jh)}{k!}.
\]
Ist $m$ fest und $f$ in einer Umgebung von $x_0$ $m$-mal stetig
differenzierbar, so gilt für $0\leq k\leq m$ bei $h\to0$
\[
 \frac{(\Delta_h^kf)(x_0)}{h^k}\longrightarrow f^{(k)}(x_0),
 \qquad
 \frac{\prod_{j=0}^{k-1}(x-x_0-jh)}{k!}
 \longrightarrow\frac{(x-x_0)^k}{k!}.
\]
Die erste Grenze ist für $k=1$ der Differenzenquotient. Allgemein
folgt sie durch wiederholte Anwendung des Hauptsatzes der
Differential- und Integralrechnung: Man mittelt
$f^{(k)}(x_0+h(t_1+\cdots+t_k))$ nacheinander über jeden Parameter
$t_i\in[0,1]$. Alle Auswertungsstellen nähern sich dabei $x_0$.
Für $k=0$ ist die Aussage unmittelbar.
Weil die Summe endlich ist, folgt für jedes feste $x$
\[
 P_{m,h}(x)\longrightarrow
 T_m(x):=\sum_{k=0}^m f^{(k)}(x_0)\frac{(x-x_0)^k}{k!}.
\]
Die Gitterstellen $x_0,x_0+h,\ldots,x_0+mh$ rücken zusammen.
Aus der Interpolation durch ihre Funktionswerte wird im Grenzfall das
Taylorpolynom, das die Ableitungen bis zur Ordnung $m$ bei $x_0$
wiedergibt. Hier bleibt $m$ fest; die Aussage ist zunächst eine
Konvergenz von Polynomen gegen $T_m$, keine Behauptung $T_m=f$.

\subsection*{Taylorpolynom, Rest und unendliche Reihe}

Für eine allgemeine Funktion ist das Taylorpolynom eine Näherung.
Die genaue Formel enthält einen Rest:
\[
 f(x)=T_m(x)+R_m(x).
\]
Ist $f$ auf einem Intervall um $x_0$ und $x$ $(m+1)$-mal stetig
differenzierbar, so gibt es für $x\ne x_0$ ein $\xi$ zwischen beiden
Stellen mit
\[
 R_m(x)=\frac{f^{(m+1)}(\xi)}{(m+1)!}(x-x_0)^{m+1}.
\]
Insbesondere ist der Rest null, wenn $f$ ein Polynom vom Grad höchstens
$m$ ist. Bei einer beliebig oft differenzierbaren Funktion stellt die
unendliche Taylorreihe $\sum_{k=0}^{\infty}
f^{(k)}(x_0)(x-x_0)^k/k!$ den Wert $f(x)$ genau dann dar, wenn
$R_m(x)\to0$. Beliebig häufige Differenzierbarkeit allein garantiert
das nicht.

Die Newtonidentität (N) hat eine andere Reichweite: Für jede Folge und
jeden natürlichen Index $n$ liefert bereits die endliche Summe bis
$n$ genau $a_n$. Eine unendliche Newtonreihe für beliebiges reelles $x$
wäre eine zusätzliche Aussage mit eigenen Konvergenzfragen.
Die Tayloranalogie erklärt die Wahl der Basis und der Koeffizienten;
für den Beweis des Folgensatzes wird keine Differentialrechnung benötigt.

\subsection*{Summation in der Newtonbasis}

Summation hebt den Basisindex um eins: Aus (S) und Teleskopieren
folgt für $N,k\in\N$
\begin{equation}\tag{A}\label{newton:sum}
 \sum_{n=0}^{N-1}\binom nk=\binom N{k+1}.
\end{equation}
Denn $\binom nk=B_{k+1}(n+1)-B_{k+1}(n)$ und $B_{k+1}(0)=0$;
für $N=0$ sind beide Seiten null. Beispielsweise ergibt die
Quadratdarstellung sofort
\[
 \sum_{n=0}^{N-1}n^2
 =\binom N2+2\binom N3
 =\frac{N(N-1)(2N-1)}6.
\]
\section{Die Kernideen der Beweistabellen}

Die Beweistabellen verwenden die kurze Induktion. Die
Operatorenschreibweise, die Interpolationsdeutung und die historischen
Quellen helfen beim Verständnis, benötigen aber keine zusätzlichen
formalen Abhängigkeiten für das Theorem.

\begin{center}
\small
\renewcommand{\arraystretch}{1.12}
\begin{tabularx}{\linewidth}{@{}p{.23\linewidth}X@{}}
\toprule
Baustein & Inhalt und Verwendung\\
\midrule
Differenzenregeln &
Linearität, $\Delta^k(\Delta a)=\Delta^{k+1}a$ und
$\Delta^r0=0$ unmittelbar aus der Rekursion, bei Bedarf durch
Induktion über $k$ beziehungsweise $r$.\\
Pascal (P) &
Fakultätsrechnung einschließlich Randfällen; fasst im
Newtonbeweis die beiden endlichen Summen zusammen.\\
Newton (N) &
Induktion über $n$ mit der Aussage \emph{für jede reelle Folge};
Anwendung auf $a$ und $\Delta a$.\\
Gradsenkung (G) &
Rekursion der Potenzdifferenzen, Linearität und Iteration;
bei exaktem Grad $r$ ist $\Delta^rp=r!c_r$.\\
Hauptschluss &
Abbruch $\Rightarrow$ feste Newtonsumme $\Rightarrow$ Polynom
$\Rightarrow$ Abbruch. Danach Eindeutigkeit durch höchste Differenz
und Bestimmung aus den ersten $d+1$ Werten.\\
\bottomrule
\end{tabularx}
\end{center}

Für ergänzende Aussagen genügt ein weiterer Baustein: Aus (B) folgt
durch Faktorisieren (S), daraus durch Iteration (K) und durch
Teleskopieren (A). Die ausgeschriebene Differenzenformel (D) ist ein
eigener kurzer Induktionsschluss; sie ist für den Hauptsatz nicht nötig.
Die Formeln (C) und die Eindeutigkeit der Interpolation folgen bereits
aus (N) und dem Abbruchkriterium.

Die Beweistabellen führen die reelle Summenrechnung, Pascal-Identität
und Differenzenregeln als eigene Hilfssätze aus. Für die Gradsenkung
verwenden sie anstelle des allgemeinen binomischen Lehrsatzes die
kürzere Rekursion
\[
 q_j(n)=(n+1)^j-n^j,\qquad
 q_{j+1}(n)=(n+1)q_j(n)+n^j.
\]
Sie liefert für $j\geq1$ induktiv den höchsten Term $j n^{j-1}$ von $q_j$.
Der anschließende Hauptbeweis ist derselbe wie oben. Die Hilfssätze
tragen den eigenen Präfix $21N$; der Hauptsatz behält seine Nummer
$21H.5.4.1$ im Leitfaden.
In den Tabellen bleiben zwei Randfälle ausdrücklich sichtbar: $n<d$ beim Ergänzen der festen Summe und $d=0$ bei Konstanten.
Der entscheidende Quantor steht in $\Delta^{d+1}a=0$:
Die Nullbedingung gilt für \emph{jeden} Index, nicht nur für die
angezeigten Stellen einer Tafel.

\section{Historischer Abriss und Originalarbeiten}

Die Darstellung verbindet zwei ältere Fragen: Wie organisiert man
wiederholte Additionen, und wie gewinnt man aus einer Wertetafel
Zwischenwerte? Das erste Problem führt zum arithmetischen Dreieck,
das zweite zur Interpolation. Die heutige Formel fasst diese Gedanken
in einer einheitlichen Schreibweise zusammen.

\textbf{Pascal: die Arithmetik des Dreiecks.}
Pascals 1665 veröffentlichter \emph{Traité du triangle arithmétique}
entwickelt die Bildung der Einträge durch Addition zweier Vorgänger
und ihre Summationsgesetze. Der Definitionsabschnitt sowie die
Conséquences II--III zeigen damit die Rechenstruktur hinter (P) und
(A) \cite{newton:pascal-source}. Das ist eine wichtige systematische
Darstellung, keine Behauptung, das Zahlendreieck sei erst bei Pascal
entstanden. Unsere Argumentation benötigt seine Additionsregel in
moderner Binomialschreibweise.

\textbf{Newton: Interpolation aus gegebenen Werten.}
Im fünften Lemma des dritten Buches der \emph{Principia} von 1687
beschreibt Newton die Konstruktion einer Kurve durch vorgegebene
Punkte: Der erste Fall benutzt gleiche Abstände, der zweite ungleiche
Abstände und entsprechend angepasste Differenzen. Schon das folgende
Lemma wendet das Verfahren auf die Bestimmung von Kometenpositionen
aus Beobachtungsdaten an \cite[Buch~III, Lemmata~V--VI]{newton:principia}.
Die Differenzentafel dient hier der Berechnung unbekannter Werte aus
bekannten Werten.

Newtons 1711 veröffentlichte \emph{Methodus differentialis} behandelt
die Interpolation ausführlicher, auch mit zentral angeordneten
Stützstellen. Die Propositionen I--III entwickeln Differenzen und
Interpolationsformen; die späteren Propositionen und das Scholium
verbinden sie mit Tabellenberechnung und Quadratur
\cite{newton:methodus}. Das historische Problem ist somit breiter
als unser Satz über Folgen. Die spezielle Darstellung an
$0,1,2,\ldots$ ist eine besonders einfache Form dieses Verfahrens.

\textbf{Taylor: Differenzenrechnung und der Übergang zu Ableitungen.}
In der \emph{Methodus incrementorum directa \& inversa} von 1715
leitet Taylor in Proposition VII eine Entwicklung aus der
Differenzentafel ab. Das zweite Korollar verbindet sie mit der
entsprechenden Entwicklung in Fluxionen, also in der damaligen Sprache
mit Ableitungen. Das Scholium zu Proposition VIII nennt den Abbruch
bei polynomialen Ausdrücken ausdrücklich
\cite[Pars~I, S.~21--24]{newton:taylor}.
Damit begegnen dort beide Leitgedanken dieser Lesefassung: die
Entwicklung aus Anfangsdifferenzen und ihr Ende bei Polynomen.

Die Bezeichnungen $\Delta$, $B_k$ und die Operatorenschreibweise sind
hier moderne Formulierungen. Die genannten Stellen dokumentieren
konkrete Beiträge und Anwendungen; sie begründen keine ausschließliche
Priorität Newtons für die gesamte Interpolationstheorie. Auch ist die
vorliegende Induktion eine heutige knappe Herleitung und keine
Übersetzung eines historischen Beweises.

\begin{thebibliography}{9}
\small
\setlength{\itemsep}{.65em}

\bibitem{newton:pascal-source}
Blaise Pascal: \emph{Traité du triangle arithmétique, avec quelques
autres petits traitez sur la mesme matière}.
Paris: Guillaume Desprez, 1665.
Hier: Definitionsabschnitt, S.~1--2, und Conséquences II--III.
\href{https://gallica.bnf.fr/ark:/12148/btv1b86262012}{Originaldigitalisat der Bibliothèque nationale de France}.
Eine \href{https://fr.wikisource.org/wiki/Trait%C3%A9_du_triangle_arithm%C3%A9tique}{Transkription des Textes}
ist ebenfalls zugänglich; sie ersetzt nicht die historische Druckausgabe.

\bibitem{newton:principia}
Isaac Newton: \emph{Philosophiae naturalis principia mathematica}.
London, 1687. Buch~III, Lemma~V, Druckseiten~481--483;
Lemma~VI, Druckseite~483.
\href{https://www.newtonproject.ox.ac.uk/view/texts/normalized/NATP00301}{Lateinischer Originaltext in der wissenschaftlichen Transkription des Newton Project}.
Die Seitenangaben beziehen sich auf die Ausgabe von 1687.

\bibitem{newton:methodus}
Isaac Newton: \emph{Methodus differentialis}.
In: William Jones (Hg.), \emph{Analysis per quantitatum series,
fluxiones, ac differentias: cum enumeratione linearum tertii ordinis}.
London, 1711, S.~93--101. Hier insbesondere Propositionen~I--III,
S.~93--97, sowie Propositionen~V--VI und Scholium, S.~99--101.
\href{https://www.e-rara.ch/download/pdf/2621379.pdf}{Originaldigitalisat des Abschnitts};
\href{https://doi.org/10.3931/e-rara-8934}{dauerhafter Nachweis des Gesamtwerks bei e-rara / ETH-Bibliothek}.

\bibitem{newton:taylor}
Brook Taylor: \emph{Methodus incrementorum directa \& inversa}.
London, 1715. Pars~I, Propositio~VII / Theorema~III, S.~21--23;
Corollarium~II, S.~23; Propositio~VIII, Scholium, S.~24.
\href{https://archive.org/details/UFIE003454_TO0324_PNI-2529_000000}{Originaldigitalisat der Ausgabe von 1715}.
Zum leichteren Lesen:
\href{https://www.17centurymaths.com/contents/brooktaylor/methodofincrementsparttwo.pdf}{Ian Bruces Ausgabe mit lateinischem Text und englischer Übersetzung, Teil I~B}.
Die Fundstellen beziehen sich auf die Druckseiten des Originals.

\end{thebibliography}

\end{document}
